\documentclass[conference]{IEEEtran}
\IEEEoverridecommandlockouts
\usepackage{cite}
\usepackage{amsmath,amssymb,amsfonts}
\usepackage{algorithmic}
\usepackage{graphicx}
\usepackage{textcomp}
\usepackage{xcolor}
\def\BibTeX{{\rm B\kern-.05em{\sc i\kern-.025em b}\kern-.08em
    T\kern-.1667em\lower.7ex\hbox{E}\kern-.125emX}}
\begin{document}

\title{Intrusion Detection in a Network Using Machine Learning and Neural Network Approaches}

\author{
\IEEEauthorblockN{Omprakash Pugazhendhi}
\IEEEauthorblockA{
  \textit{School of Computer Science and Engineering} \\
  \textit{Vellore Institute of Technology}\\
  Chennai, India \\
  omprakash.p2021@vitstudent.ac.in
}
}

\maketitle

\begin{abstract}
Network intrusion detection systems (NIDS) are a critical component of enterprise cybersecurity infrastructure. Traditional signature-based approaches fail to detect novel attacks, while anomaly-based systems suffer from high false positive rates. In this paper, we evaluate and compare supervised machine learning classifiers and neural network architectures for multi-class intrusion detection on the KDD Cup 99 and NSL-KDD benchmark datasets. We assess Random Forest, Support Vector Machine, XGBoost, Multilayer Perceptron (MLP), and LSTM across five attack categories: DoS, Probe, R2L, U2R, and Normal traffic. Our analysis reveals that XGBoost achieves the highest overall accuracy (99.3\%) on NSL-KDD while LSTM demonstrates superior performance on sequential flow-level features (98.7\%), particularly for DoS and Probe attacks that exhibit temporal patterns. We further evaluate the detection rate vs. false positive rate tradeoff across attack categories, finding that R2L and U2R attacks remain the most challenging for all evaluated models due to their low representation in training data.
\end{abstract}

\begin{IEEEkeywords}
intrusion detection, network security, machine learning, LSTM, XGBoost, NSL-KDD, anomaly detection
\end{IEEEkeywords}

\section{Introduction}

The increasing sophistication and volume of cyberattacks demands automated detection systems that go beyond rule-based signatures. Machine learning offers a data-driven approach that can generalize to novel attack patterns without manual rule authoring \cite{liao2013intrusion}. However, the performance gap between benchmark evaluations and real-world deployment remains significant: models trained on KDD Cup 99 frequently fail on live traffic due to dataset-specific statistical artifacts \cite{tavallaee2009detailed}.

NSL-KDD \cite{tavallaee2009detailed} addresses several limitations of KDD Cup 99, including the removal of redundant records that allow classifiers to achieve inflated accuracy by memorizing frequent records rather than learning meaningful decision boundaries. We evaluate on both datasets to understand the impact of these artifacts on model comparison.

Our contributions are:
\begin{enumerate}
\item A systematic comparison of five classifiers on both KDD Cup 99 and NSL-KDD.
\item Per-attack-category detection rate analysis, exposing the R2L/U2R detection problem.
\item Feature importance analysis identifying the most informative network traffic features for each attack type.
\item Analysis of detection rate vs. false positive rate tradeoffs relevant to operational NIDS deployment.
\end{enumerate}

\section{Related Work}

\subsection{Traditional NIDS Approaches}
Snort and Suricata implement signature-based detection but require continuous manual rule updates. Anomaly detection using statistical baselines (CUSUM, EWMA) provides limited coverage of complex attack patterns \cite{mukherjee1994network}.

\subsection{ML-Based NIDS}
Random Forest has been widely applied to intrusion detection \cite{pal2005random} due to its robustness to feature noise. SVM achieves competitive accuracy but scales poorly to large traffic volumes. Deep learning approaches including CNNs \cite{lecun2015deep} and LSTMs have shown promise for sequential traffic analysis.

\subsection{Benchmark Datasets}
The KDD Cup 99 dataset \cite{kdd1999} is widely used but contains known statistical artifacts. NSL-KDD \cite{tavallaee2009detailed} provides a cleaned version with balanced representation across attack categories.

\section{Dataset}

\subsection{KDD Cup 99}
494,021 training instances, 311,029 test instances. 41 features including basic TCP/IP features (duration, protocol, service, flag), content features (number of failed logins, root shell), and traffic features (same-host connections, SYN error rate).

\subsection{NSL-KDD}
125,973 training instances, 22,544 test instances. Same 41 features as KDD Cup 99 but with duplicate records removed and difficulty levels assigned to each record.

\subsection{Preprocessing}
Categorical features (protocol\_type, service, flag) encoded with one-hot encoding. Numerical features normalized to $[0, 1]$ using min-max normalization. Class labels mapped to five categories: Normal, DoS, Probe, R2L, U2R.

\section{Models}

\subsection{Traditional Classifiers}
\textbf{Random Forest}: 100 trees, max depth 20, min samples per leaf 2.\\
\textbf{SVM}: RBF kernel, $C=10$, $\gamma=0.01$.\\
\textbf{XGBoost}: 200 estimators, learning rate 0.1, max depth 8, subsample 0.8.

\subsection{Neural Networks}
\textbf{MLP}: 3 hidden layers (256, 128, 64 units), ReLU activations, dropout $p=0.3$.\\
\textbf{LSTM}: 2 LSTM layers (128 units each), 10-timestep sliding window over network flow sequences, dropout $p=0.2$.

\section{Results}

\begin{table}[htbp]
\caption{Accuracy and F1 on NSL-KDD Test Set}
\begin{center}
\begin{tabular}{|l|c|c|c|}
\hline
\textbf{Model} & \textbf{Accuracy (\%)} & \textbf{Macro F1} & \textbf{FPR (\%)} \\
\hline
Random Forest & 98.1 & 0.94 & 1.2 \\
SVM & 97.4 & 0.92 & 1.8 \\
XGBoost & \textbf{99.3} & \textbf{0.97} & 0.8 \\
MLP & 98.6 & 0.95 & 1.1 \\
LSTM & 98.7 & 0.96 & 0.9 \\
\hline
\end{tabular}
\label{tab:results}
\end{center}
\end{table}

\begin{table}[htbp]
\caption{Detection Rate (\%) by Attack Category (XGBoost)}
\begin{center}
\begin{tabular}{|l|c|}
\hline
\textbf{Category} & \textbf{Detection Rate (\%)} \\
\hline
DoS & 99.8 \\
Probe & 99.4 \\
Normal & 99.7 \\
R2L & 82.3 \\
U2R & 71.6 \\
\hline
\end{tabular}
\label{tab:per_class}
\end{center}
\end{table}

R2L and U2R attacks remain the hardest to detect across all models. U2R represents only 0.01\% of training instances in NSL-KDD, creating extreme class imbalance. SMOTE-based oversampling improved U2R detection from 71.6\% to 81.2\% at the cost of a 0.4\% increase in false positives.

\section{Discussion}

XGBoost outperforms all models on overall accuracy, which we attribute to its ability to learn non-linear decision boundaries through gradient boosting without the data requirements of deep models. LSTM's advantage in temporal pattern detection is evident in DoS and Probe attacks, which involve patterns across sequential connection records.

The R2L and U2R detection problem is a fundamental challenge: these attack types are rare, diverse, and statistically similar to normal traffic. Future work should explore few-shot learning approaches specifically for low-frequency attack categories.

\section{Conclusion}

XGBoost achieves 99.3\% accuracy on NSL-KDD with the lowest false positive rate (0.8\%), making it the recommended choice for high-throughput NIDS deployment. LSTM is preferred when sequential flow analysis is available. R2L and U2R attack categories require dedicated strategies beyond standard supervised classification. Code at \texttt{github.com/omprxkash/intrusion-detection-system-cyber}.

\begin{thebibliography}{00}
\bibitem{liao2013intrusion} H. Liao et al., ``Intrusion Detection System: A Comprehensive Review,'' J. Network Computer Applications, 2013.
\bibitem{tavallaee2009detailed} M. Tavallaee et al., ``A Detailed Analysis of the KDD CUP 99 Data Set,'' IEEE CISDA, 2009.
\bibitem{mukherjee1994network} B. Mukherjee et al., ``Network Intrusion Detection,'' IEEE Network, 1994.
\bibitem{pal2005random} M. Pal, ``Random Forest Classifier for Remote Sensing Classification,'' Int. J. Remote Sensing, 2005.
\bibitem{lecun2015deep} Y. LeCun, Y. Bengio, and G. Hinton, ``Deep Learning,'' Nature, 2015.
\bibitem{kdd1999} S. Hettich and S. Bay, ``KDD Cup 1999 Data,'' UCI Machine Learning Repository, 1999.
\end{thebibliography}

\end{document}
